\section{Машинный перевод документов}
\label{sec:translate}

\subsection{Постановка задачи}

Разработана подсистема пословного машинного перевода с английского языка на французский
(variant~5: область --- компьютерные науки и литературные эссе).
Подсистема реализует прямой перевод без перестановки слов и морфологической адаптации.

Основные возможности:
\begin{itemize}
	\item Пословный перевод с использованием базы данных словарных пар;
	\item POS-тегирование (определение частей речи) для английского и французского языков;
	\item Синтаксический разбор (построение деревьев составляющих);
	\item Частотный список слов с переводом;
	\item Управление словарём (добавление, поиск, удаление пар слов).
\end{itemize}

\subsection{Архитектура подсистемы}

Подсистема состоит из четырёх модулей:

\begin{enumerate}
	\item \texttt{translator/dictionary.py} "--- операции с базой данных PostgreSQL:
	      добавление, поиск, удаление пар слов;
	\item \texttt{translator/pos\_tagger.py} "--- POS-тегирование с использованием NLTK
	      (average perceptron tagger для EN, perceptron\_tagger\_fr для FR);
	\item \texttt{translator/syntax\_parser.py} "--- построение деревьев составляющих
	      с помощью RegexpParser;
	\item \texttt{translator/translation\_engine.py} "--- основной движок перевода:
	     lookup словаря, замена слов, сбор статистики.
\end{enumerate}

\subsection{Процесс перевода}

Алгоритм перевода текста:

\begin{enumerate}
	\item Входной текст токенизируется и разбивается на предложения.
	\item Каждое слово POS-тегируется NLTK-теггером (средний перцептрон).
	\item Для каждого слова выполняется поиск в базе данных словарных пар
	      (сначала по словоформе и POS-тегу, затем без POS-тега).
	\item Найденное французское слово подставляется вместо английского.
	\item Слова, не найденные в словаре, остаются на английском.
	\item Формируется статистика: количество слов, переведённых слов, покрытие.
\end{enumerate}

\subsection{POS-тегирование}

Для английского языка используется NLTK averaged perceptron tagger
(\texttt{averaged\_perceptron\_tagger\_eng}), для французского --- NLTK perceptron tagger
(\texttt{perceptron\_tagger\_fr}). При отсутствии модели используется эвристика на основе
суффиксов:

\begin{lstlisting}[caption={Эвристическое определение POS-тега для FR.},label={lst:pos_heuristic}]
def _heuristic_pos_fr(word: str) -> str:
    w = word.lower()
    if w.endswith(('tion', 'ment', 'age', 'ence', 'ance')):
        return 'NC'
    if w.endswith(('eux', 'if', 'ique', 'al')):
        return 'ADJ'
    if w.endswith(('er', 'ir', 're')):
        return 'V'
    return 'NC'
\end{lstlisting}

Описания POS-тегов хранятся в словаре \texttt{POS\_DESCRIPTIONS}
и передаются клиенту для отображения.

\subsection{Синтаксический разбор}

Для построения деревьев составляющих используется NLTK RegexpParser
с грамматическими правилами:

\begin{lstlisting}[caption={Грамматические правила для EN.},label={lst:en_grammar}]
EN_GRAMMAR = r"""
    NP: {<DT|PRP\$>?<JJ.*>*<NN.*>+}       # Noun phrase
    VP: {<MD>?<VB.*><NP|PP|CLAUSE>*}         # Verb phrase
    PP: {<IN><NP>}                            # Prepositional phrase
    CLAUSE: {<NP><VP>}                        # Clause
"""
\end{lstlisting}

Деревья сериализуются в JSON (вложенные dict) для передачи клиенту.
Строковое представление генерируется методом \texttt{str(tree)} NLTK.

\subsection{Словарь переводчика}

База данных PostgreSQL хранит таблицу \texttt{translation\_dictionary}:

\begin{lstlisting}[caption={Схема таблицы словаря.},label={lst:dict_schema}]
CREATE TABLE translation_dictionary (
    id SERIAL PRIMARY KEY,
    source_word TEXT NOT NULL,
    target_word TEXT NOT NULL,
    source_pos TEXT DEFAULT '',
    target_pos TEXT DEFAULT '',
    frequency INTEGER DEFAULT 1,
    UNIQUE (source_word, target_word, source_pos)
);
CREATE INDEX idx_dict_source ON translation_dictionary (LOWER(source_word));
\end{lstlisting}

Начальный словарь загружается из двух источников:
\begin{enumerate}
	\item 800+ заранее подготовленных пар слов (артикли, предлоги, местоимения,
	      глаголы, существительные, прилагательные, наречия, термины ИТ);
	\item Частотное выравнивание пар из параллельных корпусов EN/FR в проекте.
\end{enumerate}

\subsection{API методы}

\begin{lstlisting}[caption={HTTP эндпоинты подсистемы перевода.},label={lst:translate_api}]
POST /api/translate             --- перевод текста
POST /api/translate/pos         --- POS-тегирование
POST /api/translate/parse       --- синтаксический разбор
GET  /api/translate/dict        --- список словаря
POST /api/translate/dict        --- добавление пары
GET  /api/translate/dict/search --- поиск по префиксу
POST /api/translate/dict/bulk   --- массовая загрузка
POST /api/translate/dict/bootstrap --- загрузка базового словаря
POST /api/translate/dict/delete --- удаление пары
POST /api/translate/freq        --- частотный список
\end{lstlisting}

\subsection{Интерфейс пользователя}

Фронтенд (\texttt{translate.html}) содержит три вкладки:

\begin{enumerate}
	\item \textbf{Word List} "--- ввод английского текста, получение перевода
	      со статистикой (количество слов, переведённых слов, покрытие)
	      и подробной таблицей с POS-тегами и статусом перевода.
	\item \textbf{Syntax Tree} "--- ввод одного или нескольких предложений,
	      отображение деревьев составляющих (NP, VP, PP, CLAUSE).
	\item \textbf{Dictionary} "--- просмотр, поиск, добавление и удаление
	      пар слов в базе данных; кнопка ``Bootstrap Common Words''
	      для загрузки базового словаря.
\end{enumerate}

\subsection{Тестовые результаты}

Тестовый запуск перевода: \texttt{``The cat sat on the mat''}.

\begin{lstlisting}[caption={Результат перевода.},label={lst:translate_result}]
{
  "translated_text": "le cat sat sur le mat",
  "stats": {
    "total_words": 6,
    "translated": 3,
    "coverage": 50.0
  }
}
\end{lstlisting}

Слова ``The'', ``on'', ``the'' переведены (le, sur, le).
Слова ``cat'', ``sat'', ``mat'' не найдены в словаре.
После загрузки базового словаря покрытие типичного текста составляет 60--70\%.

\subsection{Выводы}

Разработана подсистема пословного машинного перевода с интеграцией в IR-систему.
Подсистема демонстрирует базовую функциональность: POS-тегирование, синтаксический
разбор, управление словарём. Ограничения: отсутствие морфологической адаптации
склонений и спряжений, отсутствие перестановки слов. Для повышения качества перевода
необходимо расширение словаря и добавление морфологической обработки.
